\subsection*{Solutions for Chapter 2}

\MCsolution{maps:ex-stability}
A contracted rational component with no special points violates the requirement of three special points. Its canonical polarization has degree $-2$. A constant map from a smooth genus-two curve is stable because its only component has genus at least two; its polarization has degree $2g-2=2$. The pullback of $A$ under a constant map has degree zero in both cases.

\medskip
\MCsolution{maps:ex-arrow}
The identity of $\PP^1$ does not give this arrow: it would require $h\circ\operatorname{id}=f$, whereas the two displayed parametrizations are different maps. The coordinate swap is the required arrow. An arrow over a noninvertible base map cannot be invertible in the total category, since any inverse would project to an inverse of the base map. Its comparison with the base-changed target is nevertheless an isomorphism.

\medskip
\MCsolution{maps:ex-cartesian}
The total-space map must lie over $w$, preserve the marked sections, preserve the map to $X$, and induce an isomorphism with the base-changed target family. For the last target equation, $f_T=f\operatorname{pr}_C$ and $\operatorname{pr}_C\widetilde b=b$, so $f_T\widetilde b=fb=e$. The product description $E\simeq C\times_S U\simeq(C\times_S T)\times_T U$ proves the cartesian-square condition.

\medskip
\MCsolution{maps:ex-isom}
The compatible source isomorphisms glue to a unique isomorphism of the source schemes over $S$, and to a pointed-curve isomorphism if the markings were respected locally. Without $h_i a_i=f_i$ there is no reason for $ha=f$. The result would therefore be an arrow in a different moduli problem. For example, the identity of a source $\PP^1$ does not identify two different parametrizations of a line as stable maps unless the maps themselves agree.

\medskip
\MCsolution{maps:ex-descent}
Projectivity supplies the very ample $A$ and hence the relatively ample $\omega(\sum\sigma_i)\otimes f^*A^3$. The cocycle equation makes the curve and line-bundle identifications consistent on triple overlaps, so that the polarized descent theorem applies. Compatibility with $f_i$ identifies the $f_i^*A^3$ factors and later lets the maps to $X$ themselves glue. Underlying pointed curves need not be stable: an unmarked degree-one $\PP^1$ is already a counterexample. It is stable as a map because its nonconstant map contributes positivity and removes source automorphisms preserving the map.

\medskip
\MCsolution{maps:ex-tree}
The dual graph is a tree with exactly one vertex carrying positive degree. If it has at least two vertices, at least one leaf has degree zero. That contracted leaf has a single node and no markings, contradicting the three-special-point requirement for a contracted rational component. With markings, a contracted leaf might have one node and two markings, and would then be stable. Thus the absence of markings is an essential hypothesis in the argument.
